% Organisation du code et des calculs.

\section{Mise en œuvre}
\label{sec-mise-en-oeuvre}

Les expériences des sections suivantes demandent plus d'une centaine d'entraînements, dont certains durent près
d'une heure. Nous avons donc organisé le code et les calculs pour pouvoir corriger, relancer et partager le travail
sans rien recalculer inutilement.

\subsection{Un code modulaire}

Tout le code réutilisable est regroupé dans un paquet Python (\texttt{src/heston\_nn}), découpé par rôle :
chargement et découpage des données, construction du réseau, entraînement, métriques, calibration et figures.
Les notebooks restent courts : ils appellent ces fonctions et présentent les résultats. Les règles communes
(grille, plages des paramètres, découpage, budget d'entraînement, graines) sont écrites une seule fois dans un
fichier de configuration (\texttt{configs/base.yaml}) ; chaque expérience n'a qu'un petit fichier qui précise ce
qui change (par exemple l'activation).

Cette organisation nous a permis de corriger et d'améliorer le code tout au long du projet sans toucher aux
notebooks : une correction dans une fonction profite à toutes les étapes. Par exemple, la gestion des points
manquants (NaN) dans la perte a été ajoutée une seule fois, dans la fonction d'entraînement, et s'applique à tous
les réseaux. Une centaine de tests automatiques vérifient chaque fonction sur des données de petite taille, ce qui
a permis de repérer les erreurs avant de lancer les longs calculs.

\subsection{Des calculs partagés et enregistrés}

Pour aller plus vite, les entraînements sont répartis entre nos 4 machines. Chaque membre de l'équipe a un numéro
de machine (de 0 à 3) : la liste de toutes les expériences, c'est-à-dire chaque configuration avec chaque graine,
est partagée de façon que chaque machine en prenne une sur quatre. Il suffit de changer ce numéro dans le notebook
pour que chacun calcule sa part.

Chaque entraînement terminé est enregistré dans son propre dossier, un par configuration et par graine :
les erreurs, l'historique d'apprentissage et le réseau. Avant de lancer un entraînement, le code vérifie s'il est
déjà fait ; si oui, il relit simplement le résultat. Un calcul n'est donc fait qu'une seule fois : si une machine
s'arrête en cours de route, il suffit de relancer la même cellule pour reprendre là où elle en était. Les résultats
de chaque machine sont ensuite rassemblés sur GitHub, et le notebook les relit tous pour tracer les courbes et
faire les tableaux. Une même référence n'est jamais recalculée : le réseau retenu à l'étape 1 sert directement de
référence à l'étape 2.

\subsection{Reproductibilité}

Toutes les sources de hasard sont fixées par une graine : tirage des paramètres, découpage des données,
initialisation des poids et ordre des lots. Avec la même configuration et la même graine, un entraînement redonne
les mêmes résultats. Le code, les configurations, les données d'entraînement et les résultats sont versionnés sur
GitHub, de sorte que chaque chiffre de ce rapport peut être retrouvé et recalculé.
